\documentclass{article}

\usepackage[margin=1in]{geometry}
\usepackage{amsmath,amsthm,amssymb}
\usepackage[spanish]{babel}
\usepackage{enumitem}

\theoremstyle{definition}
\newtheorem{definition}{Definición}[section]

% Number sets
\newcommand{\R}{\mathbb{R}}
\newcommand{\Z}{\mathbb{Z}}
\newcommand{\N}{\mathbb{N}}
\newcommand{\Q}{\mathbb{Q}}
\newcommand{\C}{\mathbb{C}}
\newcommand{\K}{\mathbb{K}}

% Functional analysis operators
\newcommand{\norm}[1]{\left\lVert #1 \right\rVert}       % Norm
\newcommand{\seminorm}[1]{\left\lVert #1 \right\rVert}   % Seminorm (same symbol)
\newcommand{\cl}[1]{\overline{#1}}                        % Closure
\newcommand{\Int}[1]{{#1}^\circ}                          % Interior
\newcommand{\bola}[2]{B(#1;#2)}                           % Ball centered at #1 with radius #2
\newcommand{\EN}[1][E]{(#1,\norm{\cdot})}                 % Normed space (E, ||.||); use \EN or \EN[F]

% Useful shortcuts
\newcommand{\sii}{\iff}                          % Si y sólo si
\newcommand{\imp}{\implies}                      % Implies
\newcommand{\setdiff}{\smallsetminus}            % Set difference
\newcommand{\diam}{\text{diam}}

% Exercise environment
\newenvironment{ejercicio}[1]{\begin{trivlist}
\item[\hskip \labelsep {\bfseries Ejercicio}\hskip \labelsep {\bfseries #1.}]}{\end{trivlist}}

\setlist[enumerate,1]{label=\alph*), leftmargin=*, itemsep=2pt}
\setlist[enumerate,2]{label=\roman*., leftmargin=*, itemsep=2pt}

\begin{document}

\title{Análisis Funcional: Práctica I}
\author{Bustos Jordi\\Espacios normados}

\maketitle

\begin{ejercicio}{1}
    Sea \( \EN \) un \( \K \)-espacio vectorial normado.
    \begin{enumerate}
        \item Demostrar que la función \( d:E\times E\to\R \) dada por
              \[
                  d(x,y):=\norm{x-y} \quad \text{para todo } x,y\in E,
              \]
              es una métrica.
        \item Probar que las siguientes funciones son continuas con la topología asociada a \(d\), la cual hace de \(E\) un Espacio Vectorial Topológico.
              \begin{enumerate}
                  \item Las traslaciones \(T_x:E\to E\) dadas por \(T_x(y)=x+y\) para \(x,y\in E\) (\(x\) fijo).
                  \item La multiplicación por un escalar \(M_x:\K\to E\) dada por \(M_x(\alpha)=\alpha x\) para \(x\in E\) y \( \alpha \in \K \) (\(x\) fijo).
                  \item La norma, \(\norm{\cdot}:E\to\R^+\) dada por \(x\mapsto \norm{x}\). Mostrar que además vale la siguiente desigualdad:
                        \[
                            \big|\norm{x}-\norm{y}\big|\leq\norm{x-y}=d(x,y) \quad \text{para todo } x,y\in E.
                        \]
              \end{enumerate}
    \end{enumerate}
\end{ejercicio}
\begin{proof}
    Sea \( \EN \) un \( \mathbb{K} \)-espacio vectorial normado.
    \begin{enumerate}
        \item Es claro que \( d(x, x) = \norm{x-x} = \norm{0} = 0 \) para todo \( x \in E \). Además \[
                  d(x, y) = \norm{x-y} = \norm{y-x} = d(y, x) \quad \text{para todo } x, y \in E.
              \]
              La desigualdad triangular se cumple porque vale para la norma y, por lo tanto, \( d \) es una métrica.
        \item Para demostrar la continuidad de las traslaciones, sea \( x \in E \) y consideremos \( T_x(y) = x + y \). Dado \( \epsilon > 0 \), si \( d(y_1, y_2) < \epsilon \), entonces
              \[
                  d(T_x(y_1), T_x(y_2)) = d(x + y_1, x + y_2) = \norm{(x + y_1) - (x + y_2)} = \norm{y_1 - y_2} = d(y_1, y_2) < \epsilon,
              \]
              lo que prueba la continuidad de \( T_x \).

              Para la multiplicación por un escalar, sea \( x \in E \) y consideremos \( M_x(\alpha) = \alpha x \). Dado \( \epsilon > 0 \), si \( |\alpha_1 - \alpha_2| < \frac{\epsilon}{\norm{x}} \), entonces
              \[
                  d(M_x(\alpha_1), M_x(\alpha_2)) = d(\alpha_1 x, \alpha_2 x) = \norm{\alpha_1 x - \alpha_2 x} = |\alpha_1 - \alpha_2| \norm{x} < \epsilon,
              \]
              lo que prueba la continuidad de \( M_x \).

              La continuidad de la norma se sigue de la desigualdad
              \[
                  \big|\norm{x}-\norm{y}\big|\leq\norm{x-y}=d(x,y) \quad \text{para todo } x,y\in E.
              \]
    \end{enumerate}
\end{proof}

\vspace{0.5cm}

\begin{ejercicio}{2}
    Sean \(E\) un espacio vectorial y \(\tilde{d}\) una métrica que verifica
    \begin{itemize}
        \item \(\tilde{d}(\lambda x,\lambda y)=|\lambda|\tilde{d}(x,y)\) para todo \(x,y\in E\) (homogeneidad).
        \item \(\tilde{d}(x+v,y+v)=\tilde{d}(x,y)\) para todo \(x,y,v\in E\) (invarianza por traslaciones).
    \end{itemize}
    Demostrar que \(\tilde{d}\) proviene de una norma.
\end{ejercicio}
\begin{proof}
    Definamos la siguiente norma \[
        \norm{x} := \tilde{d}(x, 0)
    \]
    Es fácil verificar que esta definición cumple con las propiedades de una norma:
    \begin{itemize}
        \item \(\norm{x} = 0 \iff \tilde{d}(x, 0) = 0 \iff x = 0\).
        \item \(\norm{\lambda x} = \tilde{d}(\lambda x, 0) = |\lambda| \tilde{d}(x, 0) = |\lambda| \norm{x}\).
        \item \(\norm{x+y} = \tilde{d}(x+y, 0) = \tilde{d}(x, -y) = \tilde{d}(x, 0 + (-y)) \leq \tilde{d}(x, 0) + \tilde{d}(-y, 0) = \norm{x} + \norm{y}\).
    \end{itemize}
    Por lo tanto, \(\norm{\cdot}\) es una norma y \(\tilde{d}(x,y) = \norm{x-y}\) para todo \(x,y \in E\).
\end{proof}

\vspace{0.5cm}

\begin{ejercicio}{3}
    Dar un ejemplo de una distancia que no proviene de una norma.
\end{ejercicio}
\begin{proof}
    Un ejemplo clásico de una distancia que no proviene de una norma es la distancia discreta en un conjunto \(X\) con al menos dos elementos:
    \[
        d(x,y) = \begin{cases}
            0 & \text{si } x = y,    \\
            1 & \text{si } x \neq y.
        \end{cases}
    \]
    Esta distancia no puede provenir de una norma, ya que si existiera una norma \(\norm{\cdot}\) tal que \(d(x,y) = \norm{x-y}\), entonces para cualquier \(x \neq y\) tendríamos \(\norm{x-y} = 1\). Sin embargo, para cualquier \(x \neq y\), \(\norm{2(x-y)} = 2\norm{x-y} = 2\), lo cual contradice la definición de la distancia discreta.
\end{proof}

\vspace{0.5cm}

\begin{ejercicio}{4}
    Sea \(E\) un espacio normado. Probar que si \( {\{x_n\}}_{n\in\N}\) es una sucesión de Cauchy de elementos no nulos de \(E\) que no converge a cero, entonces \( {\left \{ \dfrac{x_n}{\norm{x_n}} \right \}}_{n\in\N}\) es también una sucesión de Cauchy.
\end{ejercicio}
\begin{proof}
    Sea \( \epsilon > 0 \), \( \exists n_0 \in \N \) tal que \[
        \norm{x_n - x_m} < \epsilon \quad \forall n,m \geq n_0
    \]
    Luego, \begin{align*}
        \norm{ \dfrac{x_n}{\norm{x_n}} - \dfrac{x_m}{\norm{x_m}} } & = \norm{ \dfrac{x_n}{\norm{x_n}} - \dfrac{x_m}{\norm{x_m}} + \dfrac{x_n}{\norm{x_m}} - \dfrac{x_n}{\norm{x_m}} }                           \\
                                                                   & \leq \norm{ \dfrac{x_n}{\norm{x_n}} - \dfrac{x_n}{\norm{x_m}}} + \norm{x_n - x_m} \cdot \dfrac{1}{\norm{x_m}}                              \\
                                                                   & = \norm{ \dfrac{x_n \cdot \norm{x_m} - x_n \cdot \norm{x_n}}{\norm{x_n} \cdot \norm{x_m}} } + \norm{x_n - x_m} \cdot \dfrac{1}{\norm{x_m}} \\
                                                                   & = \dfrac{1}{\norm{x_n}\norm{x_m}} \cdot \norm{x_n \cdot (\norm{x_n} - \norm{x_m})} + \norm{x_n - x_m} \cdot \dfrac{1}{\norm{x_m}}          \\
                                                                   & = \dfrac{1}{\norm{x_n}\norm{x_m}} \cdot \norm{x_n} \cdot |\norm{x_n} - \norm{x_m}| + \norm{x_n - x_m} \cdot \dfrac{1}{\norm{x_m}}          \\
                                                                   & = \dfrac{|\norm{x_n} - \norm{x_m}|}{\norm{x_m}} + \dfrac{\norm{x_n - x_m}}{\norm{x_m}}                                                     \\
                                                                   & \leq \dfrac{\norm{x_n - x_m}}{\norm{x_m}} + \dfrac{\norm{x_n - x_m}}{\norm{x_m}}                                                           \\
                                                                   & = 2 \cdot \dfrac{\norm{x_n - x_m}}{\norm{x_m}}                                                                                             \\
                                                                   & \leq 2 \cdot \dfrac{\epsilon}{\norm{x_m}}.
    \end{align*}
    Además, como \( \norm{x_n} \not\to 0 \), existe \( c > 0 \) tal que \( \norm{x_n} \geq c \) para todo \( n \) suficientemente grande. Por lo tanto, para \( n,m \) suficientemente grandes,
    \[
        \norm{ \dfrac{x_n}{\norm{x_n}} - \dfrac{x_m}{\norm{x_m}} } \leq 2 \cdot \dfrac{\epsilon}{\norm{x_m}} \leq 2 \cdot \dfrac{\epsilon}{c}.
    \]
    Como \( \epsilon > 0 \) era arbitrario, se concluye que \( {\left \{ \dfrac{x_n}{\norm{x_n}} \right \}}_{n\in\N} \) es una sucesión de Cauchy.
\end{proof}

\vspace{0.5cm}

\begin{ejercicio}{5}
    Probar que cada uno de los siguientes espacios es normado sobre \( \R \) y decidir si es un espacio de Banach.
    \begin{enumerate}
        \item \(\ell^1:=\left \{ x={\{x_n\}}_{n\in\N}\subset\R \text{ tales que } \sum_{n\in\N}|x_n|<\infty\right \} \) con \(\norm{x}_1=\sum_{n\in\N}|x_n|\).
        \item \(\ell^\infty:=\left \{ x={\{x_n\}}_{n\in\N}\subset\R \text{ tales que } \sup_{n\in\N}|x_n|<\infty\right \} \) con \(\norm{x}_\infty=\sup_{n\in\N}|x_n|\).
        \item \(c_0:=\{x={\{x_n\}}_{n\in\N}\subset\R \text{ que convergen a } 0\} \) con \(\norm{x}_\infty=\sup_{n\in\N}|x_n|\).
        \item \(B[a,b] = \{ f : [a,b]\to\R \text{ acotadas} \} \) con \( \norm{f}_\infty = \sup_{x\in[a,b]} |f(x)| \).
        \item \(C[a,b]\) con \(\norm{f}_\infty=\sup_{x\in[a,b]}|f(x)|\).
        \item \( C^1[a,b]:=\{f:[a,b]\to\R \text{ tales que } f \text{ es de clase } C^1\} \) con
              \[
                  \norm{f}=\sup_{x\in[a,b]}|f(x)|+\sup_{x\in[a,b]}|f'(x)|.
              \]
    \end{enumerate}
\end{ejercicio}

\begin{proof}
    \begin{enumerate}
        \item \textbf{Espacio \( \ell^1 \):}

              \textit{Es espacio normado:}
              Para \( x, y \in \ell^1 \) y \( \alpha \in \R \):
              \begin{itemize}
                  \item Positividad: \( \norm{x}_1 = \sum_{n\in\N} |x_n| \geq 0 \). Si \( \norm{x}_1 = 0 \), entonces cada término \( |x_n| = 0 \), lo que implica \( x_n = 0 \) para todo \( n \), luego \( x = 0 \).
                  \item Homogeneidad: \( \norm{\alpha x}_1 = \sum_{n\in\N} |\alpha x_n| = |\alpha| \sum_{n\in\N} |x_n| = |\alpha| \norm{x}_1 \).
                  \item Desigualdad triangular: \( \norm{x+y}_1 = \sum_{n\in\N} |x_n + y_n| \leq \sum_{n\in\N} (|x_n| + |y_n|) = \norm{x}_1 + \norm{y}_1 \).
              \end{itemize}

              \textit{Es de Banach:}
              Sea \( {\{x_n\}}_{n\in\N} \) una sucesión de Cauchy en \( \ell^1 \), donde \( x_n = (x_n^{(1)}, x_n^{(2)}, \dots) \).
              Para todo \( \epsilon>0 \), existe \( N\in\N \) tal que para todo \( m,n\geq N \),
              \[
                  \norm{x_n - x_m}_1 = \sum_{k\in\N} |x_n^{(k)} - x_m^{(k)}| < \epsilon.
              \]
              En particular, para cada \( k\in\N \) fijo, \( |x_n^{(k)} - x_m^{(k)}| < \epsilon \), por lo que \( {\{x_n^{(k)}\}}_{n\in\N} \) es de Cauchy en \( \R \) y converge a un límite \( x^{(k)} \in \R \). Sea \( x := (x^{(1)}, x^{(2)}, \dots) \).

              Dado \( \epsilon>0 \), tomamos \( N \) de la condición de Cauchy. Para cualquier entero \( M \geq 1 \) y \( m,n \geq N \):
              \[
                  \sum_{k=1}^M |x_n^{(k)} - x_m^{(k)}| \leq \sum_{k=1}^\infty |x_n^{(k)} - x_m^{(k)}| < \epsilon.
              \]
              Tomando el límite cuando \( m\to\infty \), obtenemos:
              \[
                  \sum_{k=1}^M |x_n^{(k)} - x^{(k)}| \leq \epsilon.
              \]
              Como esto vale para cualquier \( M \), tomamos el supremo (o límite cuando \( M\to\infty \)) para obtener \( \sum_{k=1}^\infty |x_n^{(k)} - x^{(k)}| \leq \epsilon \). Esto prueba que \( x_n - x \in \ell^1 \) (y por ende \( x \in \ell^1 \)) y que \( x_n \to x \) en la norma \( \norm{\cdot}_1 \). Luego, \( \ell^1 \) es de Banach.

        \item \textbf{Espacio \( \ell^\infty \):}

              \textit{Es espacio normado:}
              Para \( x, y \in \ell^\infty \) y \( \alpha \in \R \):
              \begin{itemize}
                  \item Positividad: \( \norm{x}_\infty = \sup_{n\in\N} |x_n| \geq 0 \). Si \( \norm{x}_\infty = 0 \), entonces \( \sup |x_n| = 0 \implies x_n = 0 \) para todo \( n \implies x = 0 \).
                  \item Homogeneidad: \( \norm{\alpha x}_\infty = \sup_{n\in\N} |\alpha x_n| = |\alpha| \sup_{n\in\N} |x_n| = |\alpha| \norm{x}_\infty \).
                  \item Desigualdad triangular: \( \norm{x+y}_\infty = \sup_{n\in\N} |x_n + y_n| \leq \sup_{n\in\N} (|x_n| + |y_n|) \leq \sup_{n\in\N} |x_n| + \sup_{n\in\N} |y_n| = \norm{x}_\infty + \norm{y}_\infty \).
              \end{itemize}

              \textit{Es de Banach:}
              Sea \( {\{x_n\}}_{n\in\N} \) sucesión de Cauchy en \( \ell^\infty \). Para cada coordenada \( k \), \( {\{x_n^{(k)}\}}_{n\in\N} \) es de Cauchy en \( \R \) y converge a \( x^{(k)} \).
              Dado \( \epsilon>0 \), existe \( N \) tal que para \( n,m \geq N \), \( |x_n^{(k)} - x_m^{(k)}| < \epsilon \) para todo \( k \). Tomando límite con \( m \to \infty \) se obtiene \( |x_n^{(k)} - x^{(k)}| \leq \epsilon \) para todo \( k \), lo que implica \( \sup_{k} |x_n^{(k)} - x^{(k)}| \leq \epsilon \). Esto prueba que \( x \in \ell^\infty \) y que hay convergencia en norma.

        \item \textbf{Espacio \( c_0 \):}

              \textit{Es espacio normado:} Las propiedades de norma se heredan directamente de \( \ell^\infty \), al ser \( c_0 \subset \ell^\infty \).

              \textit{Es de Banach:} Basta notar que \( c_0 \) es un subespacio cerrado de \( \ell^\infty \). Si \( x_n \to x \) en \( \ell^\infty \) y cada \( x_n \in c_0 \), para cada \( \epsilon>0 \) existe \( N \) tal que \( \norm{x_N - x}_\infty < \epsilon/2 \). Como \( x_N \in c_0 \), existe \( M \) tal que para \( k \geq M \), \( |x_N^{(k)}| < \epsilon/2 \). Luego \( |x^{(k)}| \leq |x^{(k)} - x_N^{(k)}| + |x_N^{(k)}| < \epsilon \), mostrando que \( x \in c_0 \). Todo subespacio cerrado de un Banach es Banach.

        \item \textbf{Espacio \( B[a,b] \):}

              \textit{Es espacio normado:}
              \begin{itemize}
                  \item Positividad: \( \norm{f}_\infty = \sup_{x\in[a,b]} |f(x)| \geq 0 \). Si es \( 0 \), \( f(x)=0 \) para todo \( x \implies f=0 \).
                  \item Homogeneidad: \( \norm{\alpha f}_\infty = \sup_{x\in[a,b]} |\alpha f(x)| = |\alpha| \norm{f}_\infty \).
                  \item Desigualdad triangular: Idéntica a \( \ell^\infty \), usando el supremo sobre \( [a,b] \).
              \end{itemize}

              \textit{Es de Banach:}
              Si \( {\{f_n\}} \) es de Cauchy en \( B[a,b] \), converge puntualmente a una función \( f(x) \). Por el mismo argumento que en \( \ell^\infty \), la convergencia es uniforme (\( \norm{f_n - f}_\infty \to 0 \)). Al ser cada \( f_n \) acotada y converger uniformemente, \( f \) resulta acotada, por lo que \( B[a,b] \) es completo.

        \item \textbf{Espacio \( C[a,b] \):}

              \textit{Es espacio normado:} Se hereda al ser un subespacio de \( B[a,b] \) (toda función continua en un compacto es acotada).

              \textit{Es de Banach:} \( C[a,b] \) es cerrado en \( B[a,b] \). El límite uniforme de una sucesión de funciones continuas es continuo, por lo que si \( f_n \to f \) en norma infinito con \( f_n \in C[a,b] \), entonces \( f \in C[a,b] \). Por lo tanto, es completo.

        \item \textbf{Espacio \( C^1[a,b] \):}

              \textit{Es espacio normado:}
              \begin{itemize}
                  \item Positividad: \( \norm{f} = \norm{f}_\infty + \norm{f'}_\infty \geq 0 \). Si \( \norm{f} = 0 \), entonces en particular \( \norm{f}_\infty = 0 \implies f=0 \).
                  \item Homogeneidad: \( \norm{\alpha f} = \norm{\alpha f}_\infty + \norm{\alpha f'}_\infty = |\alpha|\norm{f}_\infty + |\alpha|\norm{f'}_\infty = |\alpha|\norm{f} \).
                  \item Desigualdad triangular: \( \norm{f+g} = \norm{f+g}_\infty + \norm{f'+g'}_\infty \leq \norm{f}_\infty + \norm{g}_\infty + \norm{f'}_\infty + \norm{g'}_\infty = \norm{f} + \norm{g} \).
              \end{itemize}

              \textit{Es de Banach:}
              Sea \( {\{f_n\}}_{n\in\N} \) una sucesión de Cauchy en \( C^1[a,b] \). Dado \( \epsilon>0 \), existe \( N \) tal que para \( m,n \geq N \):
              \[
                  \norm{f_n - f_m}_{C^1} = \norm{f_n - f_m}_\infty + \norm{f'_n - f'_m}_\infty < \epsilon
              \]
              Esto implica que \( \norm{f_n - f_m}_\infty < \epsilon \) y \( \norm{f'_n - f'_m}_\infty < \epsilon \).

              Por lo tanto, \( {\{f_n\}} \) y \( {\{f'_n\}} \) son sucesiones de Cauchy en el espacio de Banach \( C[a,b] \) (con la norma infinito). Existen entonces funciones \( f, g \in C[a,b] \) tales que \( f_n \to f \) y \( f'_n \to g \) uniformemente.

              Por el Teorema Fundamental del Cálculo, para cada \( x \in [a,b] \):
              \[
                  f_n(x) = f_n(a) + \int_a^x f'_n(t) dt
              \]
              Como la convergencia \( f'_n \to g \) es uniforme en \( [a,b] \), podemos intercambiar límite e integral:
              \[
                  \lim_{n\to\infty} f_n(x) = \lim_{n\to\infty} f_n(a) + \lim_{n\to\infty} \int_a^x f'_n(t) dt \implies f(x) = f(a) + \int_a^x g(t) dt
              \]
              Al ser \( g \) continua, \( f \) es derivable con derivada \( f' = g \). Como \( g \in C[a,b] \), concluimos que \( f \in C^1[a,b] \).
              Finalmente, \( \norm{f_n - f}_{C^1} = \norm{f_n - f}_\infty + \norm{f'_n - g}_\infty \to 0  \therefore \) el espacio es completo.
    \end{enumerate}
\end{proof}

\vspace{0.5cm}

\begin{ejercicio}{6}
    Sea \( \EN \) un espacio vectorial normado. Demostrar que son equivalentes:
    \begin{enumerate}
        \item[(a)] \(E\) es un espacio de Banach.
        \item[(b)] La bola cerrada \(B_E=\{x\in E:\norm{x}\leq 1\} \) es completa.
        \item[(c)] La cáscara \( S_E=\{x\in E:\norm{x}=1\} \) es completa.
        \item[(d)] Para toda sucesión \( {\{x_n\}}_{n\in\N}\subset E\) tal que \( \sum_{n\in\N} \norm{x_n} < \infty \), se tiene que \(\sum_{n\in\N}x_n\) converge en \(E\).
    \end{enumerate}
\end{ejercicio}
\begin{proof}
    \( a \iff d \) está demostrado en la \textbf{Proposición 1.1.13} del apunte de Demetrio. Haremos el siguiente camino para demostrar las equivalencias: \( a \iff b \) y \( b \implies c \).

    Veamos que \( a \implies b \), sea \( E \) un espacio de Banach i.e completo y normado es claro que \( B_E \) es normado. Veamos que es completo, para ello, dada una sucesión \( { \{ x_n \} }_{n\geq 1} \subset B_E \) de Cauchy, veamos que tiene límite en \( B_E \). Como \( { \{ x_n \} }_{n\geq1} \subset B_E \subseteq E \) y \( E \) es completo sabemos que \( x_n \to x \in E \), queremos ver que efectivamente \( x \in B_E \), es decir, que \( \norm{x} \leq 1 \).
    Como \( \norm{x_n} \leq 1 \quad \forall n \in \N \) y la norma es una función continua, vale que \[
        \lim_{n \to \infty} \norm{x_n} = \norm{x}.
    \]
    Por lo tanto, \( \norm{x} \leq 1 \) y \( x \in B_E \), lo que demuestra que \( B_E \) es completo.

    Para la recíproca, supongamos que \( B_E \) es completo. Sea \( {\{ x_n \}}_{n \geq 1} \) de Cauchy en \( E \), supongamos sin pérdida de generalidad que no tiene términos nulos y \( x_n \not \to 0 \). Por el \textbf{Ejercicio 4} la sucesión \( {\left \{ \frac{x_n}{\norm{x_n}} \right \}}_{n\geq1} \) es de Cauchy en \( B_E \) y, por lo tanto, \( \frac{x_n}{\norm{x_n}} \to L \in B_E \).
    Como \( x_n \) es de Cauchy, \( \norm{x_n} \) también lo es pues es una sucesión de números reales, por lo tanto \( \norm{x_n} \to \ell \) para algún \( \ell \geq 0 \). Así, \( x_n = \norm{x_n} \frac{x_n}{\norm{x_n}} \to \ell L \in E \), lo que demuestra que \( E \) es completo.

    Que \( b \implies c \) sale de que \( S_E \) es un subespacio cerrado de \( B_E \) y, por lo tanto, completo.
\end{proof}

\vspace{0.5cm}

\begin{ejercicio}{7}
    Probar que si \( \EN \) es un espacio vectorial normado y \(F\) es un subespacio de \(E\), entonces \(\cl{F}\) también lo es.
\end{ejercicio}
\begin{proof}
    Para demostrar que \( \cl{F} \) es subespacio debemos ver que: \begin{enumerate}
        \item \( 0 \in \cl{F} \).
        \item Si \( x,y \in \cl{F} \), entonces \( x+y \in \cl{F} \).
        \item Si \( \alpha \in \R \) y \( x \in \cl{F} \), entonces \( \alpha x \in \cl{F} \).
    \end{enumerate}
    La primera propiedad sale trivialmente pues \( 0 \in F \subseteq \cl{F} \).

    Para ver que la suma es cerrada separemos en casos, recordemos que la clausura de un conjunto es el conjunto unión sus puntos de acumulación, por lo tanto puede ocurrir que \( x, y \in F \), \( x \in F \) e \( y \in \cl{F} \setminus F \) o \( x,y \in \cl{F} \setminus F \). Si \( x,y \in F \) es claro que \( x+y \in F \subseteq \cl{F} \), si \( x \in F \) e \( y \in \cl{F} \setminus F \), entonces \( y \) es un punto de acumulación de \( F \) y por lo tanto \( x+y \) también lo es, así que \( x+y \in \cl{F} \). Finalmente, si \( x,y \in \cl{F} \setminus F \), ambos son puntos de acumulación de \( F \) y por lo tanto \( x+y \) también lo es, así que \( x+y \in \cl{F} \). La propiedad de la multiplicación por escalares se demuestra de manera análoga considerando los casos \( x \in F \) y \( x \in \cl{F} \setminus F \).

    Por lo tanto, \( \cl{F} \) es un subespacio de \( E \).
\end{proof}

\vspace{0.5cm}

\begin{ejercicio}{8}
    Sean \(N_1(\cdot)\) y \(N_2(\cdot)\) dos normas sobre un espacio vectorial \(E\). Probar que las siguientes expresiones definen una norma sobre \(E\).
    \begin{enumerate}
        \item \(\alpha N_1(\cdot)+\beta N_2(\cdot)\) para \(\alpha,\beta>0\).
        \item \( \max \{ N_1(\cdot),N_2(\cdot)\} \).
    \end{enumerate}
    Calcular la bola unidad cerrada en \(\R^2\) de la norma \(\norm{\cdot}:=\dfrac{1}{2}\norm{\cdot}_1+\dfrac{1}{2}\norm{\cdot}_\infty \). Graficar.
\end{ejercicio}
\begin{proof}
    Veamos que \(\alpha N_1(\cdot)+\beta N_2(\cdot)\) define una norma. Sean \(x,y \in E\) y \( \lambda \in \R \). Entonces:
    \begin{enumerate}
        \item \((\alpha N_1+\beta N_2)(x) = \alpha N_1(x) + \beta N_2(x) \geq 0\), y es cero si y sólo si \(x=0\).
        \item \((\alpha N_1+\beta N_2)(x+y) = \alpha N_1(x+y) + \beta N_2(x+y) \leq \alpha (N_1(x)+N_1(y)) + \beta (N_2(x)+N_2(y)) = (\alpha N_1+\beta N_2)(x) + (\alpha N_1+\beta N_2)(y)\).
        \item \((\alpha N_1+\beta N_2)(\lambda x) = \alpha N_1(\lambda x) + \beta N_2(\lambda x) = |\lambda| (\alpha N_1(x) + \beta N_2(x)) = |\lambda| (\alpha N_1+\beta N_2)(x)\).
    \end{enumerate}
    Luego, \(\alpha N_1+\beta N_2\) cumple con las tres propiedades de una norma, y por lo tanto es una norma sobre \(E\).

    \vspace{0.1cm}

    Veamos ahora que \( \max \{ N_1(\cdot),N_2(\cdot)\} \) define una norma. Sean \(x,y \in E\) y \( \lambda \in \R \). Entonces:
    \begin{enumerate}
        \item \(\max \{ N_1(x),N_2(x) \} \geq 0\), y es cero si y sólo si \(x=0\).
        \item \begin{align*}
                  \max \{ N_1(x+y),N_2(x+y) \} & \leq \max \{ N_1(x)+N_1(y),N_2(x)+N_2(y) \}              \\
                                               & \leq \max \{ N_1(x),N_2(x) \} + \max \{ N_1(y),N_2(y) \}
              \end{align*}
        \item \(\max \{ N_1(\lambda x),N_2(\lambda x) \} = \max \{ |\lambda| N_1(x),|\lambda| N_2(x) \} = |\lambda| \max \{ N_1(x),N_2(x) \} \).
    \end{enumerate}
    Luego, \( \max \{ N_1(\cdot),N_2(\cdot)\} \) cumple con las tres propiedades de una norma, y por lo tanto es una norma sobre \(E\).

    Finalmente, para la norma \( \norm{\cdot}:=\dfrac{1}{2}\norm{\cdot}_1+\dfrac{1}{2}\norm{\cdot}_\infty \), basta notar que es de la forma \(\alpha N_1+\beta N_2\) con \(\alpha=\beta=\frac{1}{2}\), que ya vimos que define una norma.
    La bola unidad cerrada en \(\R^2\) de esta norma se puede calcular considerando los puntos \((x,y)\) tales que \(\frac{1}{2}(|x|+|y|) + \frac{1}{2} \max \{|x|,|y|\} \leq 1\). Esto se puede simplificar a \(|x|+|y| \leq 2\) y \( \max \{|x|,|y|\} \leq 2\), y graficando se obtiene un octógono centrado en el origen con vértices en \((\pm 2,0)\), \((0,\pm 2)\) y \((\pm 1,\pm 1)\).
\end{proof}

\vspace{0.5cm}

\begin{ejercicio}{9}
    Sea \(E\) un espacio normado no trivial. Probar: \begin{enumerate}
        \item La clausura de cualquier bola abierta es la correspondiente bola cerrada.
        \item El diámetro de \(\bola{a}{r}\) es \(2r\).
        \item Si \(\bola{a}{r}\subset\bola{b}{s}\), entonces \(r\leq s\) y \(\norm{a-b}\leq s-r\).
    \end{enumerate}
\end{ejercicio}
\begin{proof}
    Sea \( E \) un EN no trivial.
    \begin{enumerate}
        \item Veamos que \( \cl{\bola{a}{r}} = \{ x \in E : \norm{x-a} \leq r \} \).
              Primero probaremos la inclusión \( \subseteq \). Sea \( x \in \cl{\bola{a}{r}} \). Entonces existe una sucesión \( {\{x_n\}}_{n\geq 1} \subset \bola{a}{r} \) tal que \( x_n \to x \) (es decir, \( \norm{x_n - x} \to 0 \)). Por la desigualdad triangular:
              \[
                  \norm{x - a} \leq \norm{x - x_n} + \norm{x_n - a} < \norm{x - x_n} + r
              \]
              Tomando el límite cuando \( n \to \infty \), se obtiene \( \norm{x - a} \leq r \). Por lo tanto, \( \cl{\bola{a}{r}} \subseteq \{ x \in E : \norm{x-a} \leq r \} \).

              Para la inclusión \( \supseteq \), sea \( x \in E \) tal que \( \norm{x-a} \leq r \). Si \( \norm{x-a} < r \), entonces trivialmente \( x \in \bola{a}{r} \subseteq \cl{\bola{a}{r}} \). Si \( \norm{x-a} = r \), definimos la sucesión \( x_n = a + \left(1 - \frac{1}{n}\right)(x-a) \) para \( n \geq 1 \). Notemos que:
              \[
                  \norm{x_n - a} = \left(1 - \frac{1}{n}\right)\norm{x-a} = \left(1 - \frac{1}{n}\right)r < r
              \]
              Luego, \( x_n \in \bola{a}{r} \) para todo \( n \). Además, \( \norm{x_n - x} = \frac{1}{n}\norm{x-a} \to 0 \) cuando \( n \to \infty \). Como construimos una sucesión en la bola abierta que converge a \( x \), concluimos que \( x \in \cl{\bola{a}{r}} \).

        \item Dados \( x, y \in \bola{a}{r} \), por desigualdad triangular tenemos:
              \[
                  \norm{x - y} \leq \norm{x - a} + \norm{a - y} < r + r = 2r
              \]
              Esto implica que \( \diam(\bola{a}{r}) = \sup_{x,y \in \bola{a}{r}} \norm{x-y} \leq 2r \).

              Para ver la otra desigualdad, como el espacio \( E \) es no trivial, existe un vector \( v \in E \) con \( v \neq 0 \). Sea \( u = \frac{v}{\norm{v}} \). Para cualquier \( 0 < t < r \), consideramos los puntos:
              \[
                  x_t = a + tu \quad \text{y} \quad y_t = a - tu
              \]
              Como \( \norm{x_t - a} = t < r \) y \( \norm{y_t - a} = t < r \), ambos puntos pertenecen a \( \bola{a}{r} \). La distancia entre ellos es \( \norm{x_t - y_t} = \norm{2tu} = 2t \).
              Como el diámetro es el supremo de las distancias, \( \diam(\bola{a}{r}) \geq 2t \) para todo \( t \in (0, r) \). Haciendo \( t \to r^- \), obtenemos \( \diam(\bola{a}{r}) \geq 2r \), de donde se concluye la igualdad.

        \item Si \( a = b \), la hipótesis es \( \bola{a}{r} \subset \bola{a}{s} \), lo cual implica directamente (tomando puntos a distancia \( t < r \)) que \( r \leq s \). Además, \( \norm{a-b} = 0 \leq s - r \).

              Si \( a \neq b \), consideramos el vector unitario \( u = \frac{a-b}{\norm{a-b}} \) que apunta desde \( b \) hacia \( a \). Para cualquier \( t \in (0, r) \), definimos el punto:
              \[
                  x_t = a + tu
              \]
              Es claro que \( \norm{x_t - a} = t < r \), luego \( x_t \in \bola{a}{r} \). Por hipótesis, \( x_t \in \bola{b}{s} \), por lo que \( \norm{x_t - b} < s \). Calculamos esta distancia:
              \begin{align*}
                  x_t - b & = (a - b) + t\frac{a-b}{\norm{a-b}}            \\
                          & = (a-b)\left( 1 + \frac{t}{\norm{a-b}} \right)
              \end{align*}
              Tomando norma a ambos lados, como el escalar en el paréntesis es positivo:
              \[
                  \norm{x_t - b} = \norm{a-b} \left( 1 + \frac{t}{\norm{a-b}} \right) = \norm{a-b} + t
              \]
              Sustituyendo esto en la condición \( \norm{x_t - b} < s \), obtenemos:
              \[
                  \norm{a-b} + t < s
              \]
              Como esta desigualdad vale para todo \( 0 < t < r \), podemos tomar el límite cuando \( t \to r^- \), lo que resulta en:
              \[
                  \norm{a-b} + r \leq s
              \]
              Reordenando los términos se obtiene la desigualdad buscada \( 0 \leq \norm{a-b} \leq s - r \).
    \end{enumerate}
\end{proof}

\vspace{0.5cm}

\begin{ejercicio}{10}
    Demostrar que en un espacio de Banach, toda sucesión encajada de bolas cerradas \({\{B_n\}}_{n\in\N}\) no vacías con radios \(r_n\to 0\), tiene un punto en común.
\end{ejercicio}
\begin{proof}
    Sea \({\{B_n\}}_{n\in\N}\) una sucesión encajada de bolas cerradas no vacías con radios \(r_n\to 0\). Es decir,
    \[
        B_1 \supset B_2 \supset \cdots \supset B_n \supset \cdots, \quad B_n = \cl{\bola{a_n}{r_n}} \neq \varnothing, \quad r_n \to 0.
    \]
    Como cada \(B_n\) es no vacía, podemos elegir un punto \(x_n \in B_n\). La sucesión \({x_n}\) es de Cauchy, ya que para \(m > n\),
    \[
        \norm{x_m - x_n} \leq \norm{x_m - a_m} + \norm{a_m - a_n} + \norm{a_n - x_n} \leq r_m + \norm{a_m - a_n} + r_n.
    \]
    Como \(r_n \to 0\), basta mostrar que \(\norm{a_m - a_n} \to 0\). Pero \(a_n \in B_n \subset B_m\) para \(m \leq n\), luego \(\norm{a_n - a_m} \leq r_m \to 0\). Por lo tanto, \({x_n}\) es de Cauchy y converge a algún punto \(x \in E\). Finalmente, como \(B_n\) son cerradas y \({x_n} \in B_n\), se tiene \(x \in B_n\) para todo \(n\), es decir, \(x\) es un punto en común.
\end{proof}

\vspace{0.5cm}

\begin{ejercicio}{11}
    En el espacio \( \ell^\infty \) definimos
    \[
        \norm{x}_Q=\limsup_{n\to\infty}|x_n|, \quad x=(x_n)\in\ell^\infty.
    \]
    \begin{enumerate}
        \item Probar que \(\norm{\cdot}_Q\) es una seminorma en \( \ell^\infty \) y \(\norm{x}_Q=0\) si y sólo si \(x\in c_0\).
        \item Probar que \(\norm{x}_Q=\norm{Qx}\) para todo \(x\in\ell^\infty \), donde \(Q:\ell^\infty\to\ell^\infty/c_0\) denota la proyección al cociente, y probar que este espacio está equipado con la norma cociente.
    \end{enumerate}
\end{ejercicio}
\begin{proof}
    Recordemos que el cociente de un espacio normado \( E \) con un subespacio cerrado \( F \) es el espacio \( E/F \) con la norma cociente definida por:
    \[
        \norm{\pi(x)} = d(x, F) = \inf_{y \in F} \norm{x - y}
    \]

    Veamos que \( \norm{\cdot}_Q \) es una seminorma. Sean \( x, y \in \ell^\infty \) y \( \alpha \in \R \) (o \( \C \)):
    \begin{enumerate}
        \item Positividad: \( \norm{x}_Q = \limsup_{n\to\infty} |x_n| \geq 0 \).
        \item Homogeneidad: \( \norm{\alpha x}_Q = \limsup_{n\to\infty} |\alpha x_n| = |\alpha| \limsup_{n\to\infty} |x_n| = |\alpha| \norm{x}_Q \).
        \item Desigualdad triangular: \( \norm{x+y}_Q = \limsup_{n\to\infty} |x_n + y_n| \leq \limsup_{n\to\infty} (|x_n| + |y_n|) \leq \limsup_{n\to\infty} |x_n| + \limsup_{n\to\infty} |y_n| = \norm{x}_Q + \norm{y}_Q \).
    \end{enumerate}
    Además, \( \norm{x}_Q = 0 \) si y sólo si \( \limsup_{n\to\infty} |x_n| = 0 \), lo que equivale a decir que \( x_n \to 0 \), es decir, \( x \in c_0 \).

    Probaremos que \( \norm{Qx} = \norm{x}_Q \) mostrando ambas desigualdades. Sea \( L = \norm{x}_Q = \limsup_{n\to\infty} |x_n| \).

    Por un lado, para cualquier sucesión \( c = (c_n) \in c_0 \), como \( c_n \to 0 \), los límites superiores coinciden:
    \[
        \limsup_{n \to \infty} |x_n - c_n| = \limsup_{n \to \infty} |x_n| = L
    \]
    Dado que el supremo de una sucesión siempre acota superiormente a su límite superior, tenemos \( \norm{x - c}_\infty \ge L \). Tomando el ínfimo sobre todo \( c \in c_0 \), obtenemos \( \norm{Qx} \ge L = \norm{x}_Q \).

    Por otro lado, construyamos un \( c \in c_0 \) específico. Definimos \( c_n \) como:
    \[
        c_n = \begin{cases}
            x_n - L \frac{x_n}{|x_n|} & \text{si } |x_n| > L   \\
            0                         & \text{si } |x_n| \le L
        \end{cases}
    \]
    Veamos que \( c \in c_0 \): dado \( \varepsilon > 0 \), por definición de límite superior, existe un \( N \in \N \) tal que para todo \( n \ge N \), \( |x_n| < L + \varepsilon \).
    Para estos \( n \), si \( |x_n| > L \) entonces \( |c_n| = |x_n| - L < \varepsilon \). Si \( |x_n| \le L \), entonces \( |c_n| = 0 < \varepsilon \). En cualquier caso, \( |c_n| < \varepsilon \) para \( n \ge N \), por lo que \( c \in c_0 \).

    Ahora, estimemos la norma infinito de \( x - c \):
    \[
        x_n - c_n = \begin{cases}
            L \frac{x_n}{|x_n|} & \text{si } |x_n| > L   \\
            x_n                 & \text{si } |x_n| \le L
        \end{cases}
    \]
    En ambos casos, \( |x_n - c_n| \le L \). Por lo tanto, \( \norm{x - c}_\infty \le L \).

    Como hemos encontrado un elemento en \( c_0 \) cuya distancia a \( x \) es \( \le L \), concluimos que \[ \norm{Qx} = \inf_{c \in c_0} \norm{x - c}_\infty \le L = \norm{x}_Q \]

    De ambas desigualdades se sigue que \( \norm{Qx} = \norm{x}_Q \).
\end{proof}

\vspace{0.5cm}

\begin{ejercicio}{12}
    Sean \(E\) un espacio normado y \(F\) un subespacio de \(E\). Probar que si \(F\) y \(E/F\) son completos, entonces \(E\) también lo es.
\end{ejercicio}
\begin{proof}
    Usaremos el criterio del \textbf{Ejercicio 6(d)}: basta ver que toda serie absolutamente convergente en \(E\) converge en \(E\).

    Sea \( {\{x_n\}}_{n\in\N} \subset E \) tal que \( \sum_{n\in\N} \norm{x_n} < \infty \). Como la norma cociente satisface
    \[
        \norm{\pi(x)} = \inf_{y\in F} \norm{x-y} \leq \norm{x-0} = \norm{x} \quad \text{para todo } x \in E,
    \]
    tenemos \( \norm{\pi(x_n)} \leq \norm{x_n} \) para todo \(n\), luego \( \sum_{n\in\N} \norm{\pi(x_n)} \leq \sum_{n\in\N} \norm{x_n} < \infty \).

    Como \(E/F\) es completo, por el Ejercicio 6(d) aplicado a \(E/F\), la serie \( \sum_{n\in\N} \pi(x_n) \) converge en \(E/F\) a algún \( \xi \in E/F \).

    Sea \( s_N := \sum_{n=1}^N x_n \) la sucesión de sumas parciales en \(E\). Como \( \pi \) es lineal, \( \pi(s_N) \to \xi \) en \(E/F\). Elijamos \( x \in E \) tal que \( \pi(x) = \xi \) (posible pues \( \pi \) es sobreyectiva). Entonces
    \[
        \norm{\pi(s_N - x)} = \norm{\pi(s_N) - \xi} \to 0, \quad \text{es decir,} \quad d(s_N - x, F) \to 0.
    \]

    Para cada \(N\), por definición de ínfimo, elijamos \( y_N \in F \) tal que
    \[
        \norm{s_N - x - y_N} < d(s_N - x, F) + \frac{1}{N} =: \epsilon_N \xrightarrow{N\to\infty} 0.
    \]

    Veamos que \( {\{y_N\}} \) es de Cauchy en \(F\). Para \( M < N \):
    \begin{align*}
        \norm{y_N - y_M} & \leq \norm{y_N - (s_N-x)} + \norm{(s_N-x)-(s_M-x)} + \norm{(s_M-x)-y_M}                            \\
                         & < \epsilon_N + \norm{s_N - s_M} + \epsilon_M = \epsilon_N + \epsilon_M + \norm{\sum_{n=M+1}^N x_n} \\
                         & \leq \epsilon_N + \epsilon_M + \sum_{n=M+1}^N \norm{x_n}.
    \end{align*}
    Como \( \epsilon_N \to 0 \) y la cola de la serie \( \sum \norm{x_n} \) tiende a \(0\), el lado derecho tiende a \(0\) cuando \( M,N \to \infty \), por lo que \( {\{y_N\}} \) es de Cauchy en \(F\).

    Como \(F\) es completo, \( y_N \to y \) para algún \( y \in F \). Finalmente,
    \[
        \norm{s_N - (x+y)} \leq \norm{s_N - x - y_N} + \norm{y_N - y} < \epsilon_N + \norm{y_N - y} \xrightarrow{N\to\infty} 0,
    \]
    por lo que \( s_N \to x+y \in E \). Esto muestra que \( \sum_{n\in\N} x_n \) converge en \(E\), y por el Ejercicio 6(d), \(E\) es completo.
\end{proof}

\vspace{0.5cm}

\begin{ejercicio}{13}
    Sean \(E\) un espacio normado y \(F\) un subespacio cerrado de \(E\). Probar:
    \begin{enumerate}
        \item Si \(E\) es separable entonces también lo es \(E/F\).
        \item Si \(F\) y \(E/F\) son separables entonces también lo es \(E\).
        \item Dar un ejemplo en el cual \(E/F\) sea separable pero \(E\) no.
    \end{enumerate}
\end{ejercicio}
\begin{proof}
    \begin{enumerate}
        \item Sea \( D = {\{d_n\}}_{n\in\N} \) un conjunto numerable denso en \(E\). Como vimos en el Ejercicio 12, la proyección al cociente satisface \( \norm{\pi(x)} \leq \norm{x} \) para todo \( x \in E \), es decir, \( \pi \) es \(1\)-Lipschitz.

              Sea \( \xi \in E/F \), digamos \( \xi = \pi(x) \) con \( x \in E \), y sea \( \epsilon > 0 \). Como \(D\) es denso, existe \( d_n \in D \) tal que \( \norm{x - d_n} < \epsilon \). Luego
              \[
                  \norm{\xi - \pi(d_n)} = \norm{\pi(x - d_n)} \leq \norm{x-d_n} < \epsilon.
              \]
              Por lo tanto \( \pi(D) = {\{\pi(d_n)\}}_{n\in\N} \) es un conjunto numerable denso en \(E/F\), es decir, \(E/F\) es separable.

        \item Sea \( D_F \) un conjunto numerable denso en \(F\) y sea \( {\{\pi(x_n)\}}_{n\in\N} \) (con \( x_n \in E \)) numerable denso en \(E/F\). Definamos
              \[
                  D := \{ x_n + f : n \in \N,\ f \in D_F \},
              \]
              que es numerable por ser unión numerable de conjuntos numerables. Veamos que \(D\) es denso en \(E\).

              Sea \( x \in E \) y \( \epsilon > 0 \). Como \( {\{\pi(x_n)\}} \) es denso en \(E/F\), existe \(n\) tal que \( \norm{\pi(x) - \pi(x_n)} = \norm{\pi(x-x_n)} = d(x - x_n, F) < \epsilon/2 \). Por definición de ínfimo, existe \( y \in F \) tal que \( \norm{(x - x_n) - y} < \epsilon/2 \). Como \(D_F\) es denso en \(F\), existe \( f \in D_F \) tal que \( \norm{y - f} < \epsilon/2 \). Luego
              \[
                  \norm{x - (x_n + f)} \leq \norm{x - x_n - y} + \norm{y - f} < \epsilon.
              \]
              Por lo tanto \(D\) es denso numerable en \(E\), es decir, \(E\) es separable.

        \item Sea \( E = \ell^\infty \), que no es separable. Por Hahn Banach, existe un funcional lineal acotado no nulo \( f: E \to \R \) (basta extender un funcional no nulo de un subespacio de dimensión \(1\)). Sea \( F := \ker f \), que es un subespacio cerrado de codimensión \(1\) (por ser \(f\) continuo y no nulo).

              Entonces \( E/F \) tiene dimensión \(1\) (es isomorfo a \( \R \) vía el mapa inducido por \(f\)), y todo espacio de dimensión finita es separable. Sin embargo, \( E = \ell^\infty \) no es separable. Este es el ejemplo buscado.
    \end{enumerate}
\end{proof}

\vspace{0.5cm}

\begin{ejercicio}{14}
    Sean \(E\) un espacio vectorial normado y \(F\subset E\) un subespacio de dimensión finita. Probar que para todo \(x\notin F\) existe \(y_0\in F\) que realiza la distancia, es decir,
    \[
        \norm{x-y_0}=\inf_{y\in F}\norm{x-y}.
    \]
\end{ejercicio}
\begin{proof}
    Fijemos \( x \notin F \) y sea \( d := \inf_{y\in F} \norm{x-y} \). Por definición de ínfimo, existe una sucesión \( {\{y_n\}}_{n\in\N} \subset F \) tal que \( \norm{x - y_n} \to d \).

    Veamos que \( {\{y_n\}} \) está acotada. Para \(n\) suficientemente grande, \( \norm{x-y_n} < d+1 \), y por la desigualdad triangular,
    \[
        \norm{y_n} \leq \norm{y_n - x} + \norm{x} < d + 1 + \norm{x} =: R.
    \]
    Por lo tanto, salvo finitos términos, \( {\{y_n\}} \subset \cl{\bola{0}{R}} \cap F \).

    Como \(F\) tiene dimensión finita \( \cl{\bola{0}{R}} \cap F \) es compacto. Luego, existe una subsucesión \( {\{y_{n_k}\}}_{k\in\N} \) que converge a algún \( y_0 \in F \). Por la continuidad de la norma,
    \[
        \norm{x - y_0} = \lim_{k\to\infty} \norm{x - y_{n_k}} = d,
    \]
    ya que \( {\{\norm{x-y_{n_k}}\}} \) es una subsucesión de \( {\{\norm{x-y_n}\}} \), que converge a \(d\). Por lo tanto, \( y_0 \) realiza la distancia de \(x\) a \(F\).
\end{proof}

\vspace{0.5cm}

\begin{ejercicio}{15}
    Dar un ejemplo de dos subespacios cerrados de un espacio de Banach tal que su suma no sea cerrada.
\end{ejercicio}
\begin{proof}
    Sea \( H := \ell^2 \oplus \ell^2 \), con la norma \( \norm{(x,y)}^2 = \norm{x}_2^2 + \norm{y}_2^2 \), que es un espacio de Banach. Definamos el operador \( T:\ell^2\to\ell^2 \) dado por
    \[
        T(x_1,x_2,x_3,\dots) := \left(x_1, \frac{x_2}{2}, \frac{x_3}{3}, \dots\right).
    \]
    Como \( \norm{Tx}_2^2 = \sum_{n\in\N} \frac{x_n^2}{n^2} \leq \sum_{n\in\N} x_n^2 = \norm{x}_2^2 \), el operador \(T\) es lineal, acotado (con \( \norm{T}\leq 1 \)) e inyectivo.

    El rango de \(T\) es \( \operatorname{Ran}(T) = \{ y=(y_n) \in \ell^2 : (n y_n)_{n} \in \ell^2 \} \), que contiene a todas las sucesiones de soporte finito y, por lo tanto, es denso en \( \ell^2 \). Sin embargo, \( \operatorname{Ran}(T) \neq \ell^2 \): por ejemplo, la sucesión \( y = (1/n)_{n\in\N} \in \ell^2 \) (pues \( \sum 1/n^2 <\infty \)) no está en \( \operatorname{Ran}(T) \), ya que requeriría \( x_n = n y_n = 1 \) para todo \(n\), y \( (1,1,1,\dots) \notin \ell^2 \). Por lo tanto, \( \operatorname{Ran}(T) \) es denso pero no cerrado en \( \ell^2 \).

    Definamos ahora los subespacios de \(H\):
    \[
        M := \ell^2 \times \{0\}, \qquad N := \{ (x, Tx) : x \in \ell^2 \} \quad (\text{el gráfico de } T).
    \]
    \(M\) es cerrado por ser un sumando directo de \(H\). \(N\) también es cerrado: si \( (x_n, Tx_n) \to (x,y) \) en \(H\), entonces \( x_n \to x \) y \( Tx_n \to y \); pero por continuidad de \(T\), \( Tx_n \to Tx \), luego \( y = Tx \) y \( (x,y) \in N \).

    Calculemos la suma:
    \[
        M + N = \{ (x,0) + (z,Tz) : x,z \in \ell^2 \} = \{ (x+z, Tz) : x,z \in \ell^2 \}.
    \]
    Para \(z\) fijo, como \(x\) recorre todo \( \ell^2 \), la primera coordenada \( x+z \) también recorre todo \( \ell^2 \). Por lo tanto,
    \[
        M+N = \ell^2 \times \operatorname{Ran}(T).
    \]
    Como \( \operatorname{Ran}(T) \) es denso pero no cerrado en \( \ell^2 \), el conjunto \( M+N = \ell^2 \times \operatorname{Ran}(T) \) es denso en \(H\) (su clausura es \( \ell^2 \times \ell^2 = H \)) pero no es cerrado, pues \( \operatorname{Ran}(T) \neq \ell^2 \).

    Así, \(M\) y \(N\) son subespacios cerrados del espacio de Banach \(H\) cuya suma \(M+N\) no es cerrada.
\end{proof}

\end{document}
